% Fixed134 substantive insert:
% exact twisted and untwisted marked Fox matrices, integral H2 trace,
% the degree-shifted Fox Bockstein, mod-four Wu class, and odd-index-3
% restriction compatibility.  This does not assert general full-group
% projective Fox completeness outside the two rank-one coefficients.

\subsection{Integral Fox--Tate trace on the three Schreier relation faces}
\label{subsec:dyadic-f134-relative-fox-tate}

The group-level seven-generator, three-relator
Schreier presentation of
Theorem~\ref{thm:dyadic-f133-schreier-pro2-presentation}
carries a previously hidden source-owned
\emph{top-cell trace}.  Its unweighted
sum over the three wild relator cells
simultaneously detects ordinary
two-torsion and the cyclotomic integral
fundamental class.  More importantly,
for the selected index-three extension
$F/\mathbf Q_2$, its interaction with
restriction is an exact multiplication
by $3$, not only an abstract
local-duality statement.
No chosen inverse of the
pro-$2$ Nielsen automorphism
is needed to read out these two
top-degree classes.

Put $P_F=G_F(2)$, with
$F=\mathbf Q_2(\sqrt[3]2)$, and retain
the three relators $R_i$ and
the ordered seven generators
\[
(s,a_0,a_1,a_2,v_0,v_1,v_2)
\]
of
\eqref{eq:dyadic-f133-seven-three-presentation}.
There are two coefficient systems:
the trivial $\mathbf Z_2$
and the integral cyclotomic
module $\mathbf Z_2(1)$.
Write $\vartheta:P_F\to\mathbf Z_2^\times$
for the canonical orientation.
By the fully marked
Roe--Turturean reciprocity
normalization \cite{RoeTurtureanGQ2},
\begin{equation}\label{eq:dyadic-f134-orientation-values}
\boxed{
\vartheta(s)=1,\qquad
\vartheta(a_i)=-1,\qquad
\vartheta(v_i)=-\frac13
\quad(i=0,1,2).
}
\end{equation}
These values come from
$\vartheta(\sigma)=\vartheta(\tau)=1$,
$\vartheta(x_0)=-1$ and
$\vartheta(x_1)=-1/3$.
They descend to the seven-generator
Schreier quotient and identify
the cyclotomic character of $G_F(2)$.

\begin{lemma}[Twisted completed Fox rule for the Schreier cube roots]
\label{lem:dyadic-f134-twisted-cubic-fox}
Let $D$ be a continuous
$\vartheta$-crossed
derivation into
$\mathbf Z_2(\vartheta)$,
so that
\[
D(gh)=D(g)+\vartheta(g)D(h).
\]
For an element $w$
whose cyclotomic
value $\vartheta(w)$
is the indicated
unit, the continuous
$1/3$ power obeys
\begin{equation}\label{eq:dyadic-f134-cuberoot-fox-factors}
\boxed{
\begin{array}{c|c|c}
\vartheta(w)&\vartheta(w^{1/3})&
D(w^{1/3})/D(w)\\ \hline
-1&-1&1\\
-1/27&-1/3&9/7
\end{array}}
\end{equation}
where the ratio denotes
the scalar multiplying
\emph{every} generator's
Fox contribution, not a
division by a possibly
zero cocycle value.
For the actual
Schreier packets,
\[
D(u_{0,i})=D(a_i a_{i+1}a_{i+2}),\qquad
D(u_{1,i})=
\frac97D(v_i v_{i+1}v_{i+2}).
\]
All coefficients are
integral in $\mathbf Z_2$.
\end{lemma}

\begin{proof}
For an integer $n\ge1$,
\[
D(w^n)=
(1+\vartheta(w)+\cdots+
\vartheta(w)^{n-1})D(w).
\]
Both sides extend
continuously to
$n\in\mathbf Z_2$.
For $n=1/3$, if
$\vartheta(w)\ne1$,
the multiplier is
\[
\frac{\vartheta(w)^{1/3}-1}
{\vartheta(w)-1}.
\]
For $\vartheta(w)=-1$
the continuous $1/3$ power
is again $-1$,
and the ratio is $1$.
For $\vartheta(w)=-1/27$
its unique $1/3$ power
in $1+4\mathbf Z_2$
is $-1/3$, giving
\[
\frac{-1/3-1}{-1/27-1}
=\frac97.
\]
The numerator/denominator
interpretation is legitimate
as a scalar identity
over $\mathbf Q_2$, and
the result belongs to
$\mathbf Z_2$.
The corresponding
packet cyclotomic values
are the products of the
three marked generator
values from
\eqref{eq:dyadic-f134-orientation-values}.
\end{proof}

\begin{theorem}[The complete three-face Fox matrices in both integral twists]
\label{thm:dyadic-f134-exact-two-fox-matrices}
Let
$\mathcal C_{\rm Sch}^\bullet(P_F,B)$
be the completed
three-relator
Fox cochain complex
\[
B\xrightarrow{d^0}B^7
\xrightarrow{d^1}B^3,
\]
for $B=\mathbf Z_2$
or $\mathbf Z_2(1)$.
In the ordered
seven-generator
coordinate system
and the ordered
relation faces $(R_0,R_1,R_2)$,
the \emph{complete}
trivial-coefficient
differentials are
\begin{equation}\label{eq:dyadic-f134-trivial-fox-matrix}
\boxed{
d_{\bf1}^0=0,\qquad
d_{\bf1}^1=
\begin{pmatrix}
0&-2/3&4/3&4/3&2/3&-1/3&-1/3\\
0&4/3&-2/3&4/3&-1/3&-1/3&2/3\\
0&4/3&4/3&-2/3&-1/3&2/3&-1/3
\end{pmatrix}.
}
\end{equation}
The \emph{complete}
cyclotomic twisted
differentials are
\begin{equation}\label{eq:dyadic-f134-twisted-fox-matrix}
\boxed{
\begin{aligned}
d_\vartheta^0(1)
&=(0,-2,-2,-2,-4/3,-4/3,-4/3)^t,\\
d_\vartheta^1
&=
\begin{pmatrix}
0&0&-2&2&6/7&-9/7&3/7\\
0&2&0&-2&3/7&27/7&-30/7\\
0&-2&2&0&-9/7&-18/7&27/7
\end{pmatrix}.
\end{aligned}}
\end{equation}
Both are actual
integral matrices
over $\mathbf Z_2$,
and
\[
d_{\bf1}^1d_{\bf1}^0
=d_\vartheta^1d_\vartheta^0=0.
\]
The twisted matrix has
the exact \emph{columnwise}
trace-zero identity
\begin{equation}\label{eq:dyadic-f134-cyclotomic-fox-zero-trace}
\boxed{\qquad
(1,1,1)d_\vartheta^1=0,
\qquad
\det\begin{pmatrix}
-9/7&3/7\\
27/7&-30/7
\end{pmatrix}
=\frac{27}{7}\in\mathbf Z_2^\times.
\qquad}
\end{equation}
The trivial matrix has
\[
(1,1,1)d_{\bf1}^1
=(0,2,2,2,0,0,0)
\]
and a rank-two
unit minor
\begin{equation}\label{eq:dyadic-f134-trivial-two-row-unit}
\det\begin{pmatrix}
-1/3&-1/3\\
-1/3&2/3
\end{pmatrix}
=-\frac13\in\mathbf Z_2^\times.
\end{equation}
These explicit
relations determine
the \emph{entire
top-coboundary images}:
\begin{equation}\label{eq:dyadic-f134-images-by-trace}
\boxed{
\begin{aligned}
\operatorname{im}d_\vartheta^1
&=\{q\in\mathbf Z_2^3:
q_0+q_1+q_2=0\},\\
\operatorname{im}d_{\bf1}^1
&=\{q\in\mathbf Z_2^3:
q_0+q_1+q_2\in2\mathbf Z_2\}.
\end{aligned}}
\end{equation}
\end{theorem}

\begin{proof}
The augmented
trivial matrix is the
three-by-seven row
\eqref{eq:dyadic-f133-augmented-three-seven}.
For the twisted
matrix, compute each
$\vartheta$-crossed
derivation of the three
explicit group words
\eqref{eq:dyadic-f133-u-and-d-words}
and
\eqref{eq:dyadic-f133-three-wild-relators}.
Use the inverse rule
\[
D(g^{-1})=-\vartheta(g)^{-1}D(g),
\]
the product rule,
and the two cube-root
factors from
Lemma~\ref{lem:dyadic-f134-twisted-cubic-fox}.
The values of the
intermediate words are
\[
\vartheta(u_{0,i})=-1,\quad
\vartheta(u_{1,i})=-1/3,\quad
\vartheta(d_i)=1,\quad
\vartheta(z_i)=-1,\quad
\vartheta(g)=1.
\]
Commutator derivative
and conjugation then give
exactly
\eqref{eq:dyadic-f134-twisted-fox-matrix};
each relation evaluates
to cyclotomic value one,
so the completed Fox
cochain formula applies.

The two displayed
$2\times2$ determinants
are odd $2$-adic units.
Thus in each complex
the first two relation
coordinates can be
used to generate the
rank-two trace-zero
hyperplane integrally.
For the twisted
matrix every column
already lies in that
hyperplane and the
unit minor proves
that the image equals
it.  For the trivial
matrix the three
$v_i$ columns generate
the trace-zero
hyperplane and any
$a_i$ column has trace
$2$.  Hence the image
is exactly the full
preimage of
$2\mathbf Z_2$
under the sum map.
This proves
\eqref{eq:dyadic-f134-images-by-trace}
and the cochain identity.
All asserted fractions
have odd denominators,
so there has been
no inversion of $2$.
\end{proof}

\begin{theorem}[Primitive top-face trace and a Fox Bockstein degree shift]
\label{thm:dyadic-f134-top-trace-bockstein-shift}
Define a single
integral linear
functional on the
three indexed
wild faces:
\begin{equation}\label{eq:dyadic-f134-primitive-top-functional}
\boxed{\qquad
\operatorname{Tr}_F:
\mathbf Z_2^3\longrightarrow\mathbf Z_2,
\qquad
(q_0,q_1,q_2)\longmapsto q_0+q_1+q_2.
\qquad}
\end{equation}
The same relation-face
functional induces
two different
top-degree isomorphisms:
\begin{equation}\label{eq:dyadic-f134-top-class-identifications}
\boxed{
\begin{aligned}
H^2(P_F,\mathbf Z_2(1))
&\xrightarrow[\sim]{\operatorname{Tr}_F}
\mathbf Z_2,\\
H^2(P_F,\mathbf Z_2)
&\xrightarrow[\sim]{\operatorname{Tr}_F\bmod2}
\mathbf Z/2\mathbf Z.
\end{aligned}}
\end{equation}
Moreover the
\emph{exact} two
integral cochain
homotopy types are
\begin{equation}\label{eq:dyadic-f134-twisted-untwisted-normal-forms}
\boxed{
\begin{aligned}
\mathcal C_{\rm Sch}^\bullet(P_F,\mathbf Z_2)
\simeq{}&
\mathbf Z_2[0]\oplus
\mathbf Z_2^4[-1]\oplus
[\mathbf Z_2\xrightarrow{\ 2\ }
 \mathbf Z_2]_{[1,2]},\\
\mathcal C_{\rm Sch}^\bullet(P_F,\mathbf Z_2(1))
\simeq{}&
[\mathbf Z_2\xrightarrow{\ 2\ }
 \mathbf Z_2]_{[0,1]}
\oplus\mathbf Z_2^4[-1]
\oplus\mathbf Z_2[-2].
\end{aligned}}
\end{equation}
Thus:
\begin{equation}\label{eq:dyadic-f134-h0h1h2-explicit}
\boxed{
\begin{array}{c|ccc}
&H^0&H^1&H^2\\ \hline
\mathbf Z_2&
\mathbf Z_2&
\mathbf Z_2^4&
\mathbf Z/2\\
\mathbf Z_2(1)&
0&
\mathbf Z_2^4\oplus\mathbf Z/2&
\mathbf Z_2 .
\end{array}}
\end{equation}
The unique two-primary
elementary pair
moves from
degrees $(1,2)$
in the untwisted
Fox complex to
degrees $(0,1)$
in the cyclotomic
one.  This is a
literal integral
Fox--Bockstein
\emph{degree shift},
not a rational
Euler-characteristic
calculation.

In the minimal
five-generator Fox
complex of
Theorem~\ref{thm:dyadic-f133-integral-schreier-fox-chain},
the top relation-face
coordinate induced by
the chosen primitive
relation product
$R_0R_1R_2$ is
precisely
$\operatorname{Tr}_F$.
Its primitiveness is
visible already
in the unit minor
of
\eqref{eq:dyadic-f134-cyclotomic-fox-zero-trace};
it does not depend
on computing the
nonlinear inverse
Nielsen word
$\alpha^{-1}$
from fixed133.
\end{theorem}

\begin{proof}
By
\eqref{eq:dyadic-f134-images-by-trace},
the quotient of
$\mathbf Z_2^3$
by the twisted
top coboundaries
is freely generated
by the sum map,
while for trivial
coefficients its
quotient is the
sum modulo two.
Since the
seven-generator
Schreier Fox
complex is
integrally chain
equivalent to the
genuine rank-five
Demu\v{s}kin
resolution by
Theorem~\ref{thm:dyadic-f133-integral-schreier-fox-chain},
these quotients
compute the actual
continuous local
$H^2$ groups.

The matrix
$d_{\bf1}^1$
has two unit
Smith factors
and one factor
of $2$;
$d_{\bf1}^0=0$,
and its integral
kernel is
free of rank four.
This gives the
first elementary
decomposition.
The twisted
$d_\vartheta^1$
has two unit
factors and
zero third row,
so its kernel
is free of rank five.
The degree-zero
map has one
nonzero Smith
factor $2$,
because
\[
d_\vartheta^0(1)
=2(0,-1,-1,-1,-2/3,-2/3,-2/3)^t
\]
and the vector
in parentheses
is primitive
over $\mathbf Z_2$.
As $d_\vartheta^1
d_\vartheta^0=0$,
this degree-zero
image belongs
to the rank-five
kernel.
Cancelling the
two degree-one--two
unit pairs leaves
the second
displayed
homotopy decomposition.
The claimed
cohomology groups
follow at once.

Under the
group-level
Tietze transformation,
the primitive
surviving relator
is the product
$R_0R_1R_2$
before the two
generator substitutions.
For any crossed
derivation, the
cocycle value
on a product
of \emph{relation
words with
cyclotomic value one}
is the sum of
its values on the
individual words.
This makes the
induced degree-two
cochain coordinate
literally
$q_0+q_1+q_2$.
The same
additivity holds
for the trivial
coefficient action.
The two primitive
relation directions
have been
removed by
integral Nielsen
pivots, so the
top sum remains
the complete
surviving face
coordinate.
\end{proof}

\begin{corollary}[The Fox--Wu class in the marked five-generator basis]
\label{cor:dyadic-f134-marked-wu}
Let
$(\chi_s,\chi_{a_0},
\chi_{a_1},\chi_{a_2},
\chi_{v_2})$
be the mod-$2$
basis dual to
the five
minimal generators
from fixed133.
The restriction of
the canonical
dyadic cyclotomic
orientation
has the explicit
mod-four Wu class
\begin{equation}\label{eq:dyadic-f134-marked-wu-class}
\boxed{\qquad
w_F=
\frac{\vartheta-1}{2}\bmod2
=\chi_{a_0}+
\chi_{a_1}+\chi_{a_2}.
\qquad}
\end{equation}
For every
$\chi\in H^1(P_F,\mathbf F_2)$
the full quadratic
Fox cup form of
\eqref{eq:dyadic-f133-cup-rank-five}
satisfies
\begin{equation}\label{eq:dyadic-f134-fox-wu-compatibility}
\boxed{\qquad
\beta_1(\chi)=
\chi\smile\chi=
\chi\smile w_F.
\qquad}
\end{equation}
This is a
source-native
compatibility
between the
marked reciprocity
orientation and
the surviving
relation's
nonalternating
Magnus symbol.
\end{corollary}

\begin{proof}
In the
seven-generator
Schreier marking,
$\vartheta(s)=1$,
$\vartheta(a_i)=-1$,
and
$\vartheta(v_i)=-1/3
\equiv1\bmod4$.
The two
Nielsen-eliminated
$v$ directions
are relator
words and have
cyclotomic value
one, so the
same character
values descend
to the five
minimal generators.
Reducing
$(\vartheta-1)/2$
modulo two gives
the displayed
sum of the
three $a_i$
dual coordinates.
The mod-two
Fox cup matrix
of
\eqref{eq:dyadic-f133-cup-rank-five}
has diagonal
ones exactly
in those
three coordinates
and has its
remaining
hyperbolic block
in $(s,v_2)$.
Therefore
$\chi\smile\chi$
equals
$\chi\smile w_F$
for every $\chi$.
The identification
of $\beta_1$
with the
degree-one cup
square is
the standard
mod-two
Bockstein identity,
also proved in
Theorem~\ref{thm:dyadic-fox-wu-bockstein}.
\end{proof}

\begin{theorem}[The three-fold top-cell restriction diagram]
\label{thm:dyadic-f134-top-transfer-three}
Let $\Gamma=G_{\mathbf Q_2}$,
$H=G_F$,
$[F:\mathbf Q_2]=3$,
and use the actual
Roe--Turturean marked
generators
$(\sigma,\tau,x_0,x_1)$.
For both rank-one
coefficient systems
$B=\mathbf Z_2$
and $B=\mathbf Z_2(1)$,
let
\[
\mathcal J_\Gamma^\bullet(B):
\quad B\to B^4\to B^2
\]
be the marked
two-face \emph{formal}
Fox coefficient
complex.
In source-generator
coordinates
$(\sigma,\tau,x_0,x_1)$
and face order
(tame,wild),
its literal
completed Fox
differentials are
\begin{equation}\label{eq:dyadic-f134-full-marked-two-twist-rows}
\boxed{
\begin{aligned}
d_{\Gamma,\bf1}^0&=0,&
d_{\Gamma,\bf1}^1&=
\begin{pmatrix}
0&-1&0&0\\
0&3&2&0
\end{pmatrix},
\\
d_{\Gamma,\vartheta}^0(1)
&=(0,0,-2,-4/3)^t,&
d_{\Gamma,\vartheta}^1&=
\begin{pmatrix}
0&-1&0&0\\
0&-3&0&0
\end{pmatrix}.
\end{aligned}}
\end{equation}
Define the global
face functionals
\begin{equation}\label{eq:dyadic-f134-global-marked-top-readout}
\boxed{
\begin{aligned}
\operatorname{Tr}_{\Gamma,\bf1}(t,w)
&=w+3t\pmod2,\\
\operatorname{Tr}_{\Gamma,\vartheta}(t,w)
&=w-3t\in\mathbf Z_2.
\end{aligned}}
\end{equation}
They annihilate
the appropriate
degree-one
Fox coboundaries,
and induce
isomorphisms from
the marked
two-face cokernels
onto
$\mathbf Z/2$
and
$\mathbf Z_2$,
respectively.

After normalizing
a global top
cochain to have
zero tame-face
coordinate, the
index-three
Schreier pullback
of its wild-face
value $c$ is the
actual triple
\begin{equation}\label{eq:dyadic-f134-three-face-restriction}
\boxed{\qquad
c\longmapsto(c,c,c)
\quad\text{on }
(R_0,R_1,R_2).
\qquad}
\end{equation}
Hence the
intrinsic local
degree-two
restriction
satisfies, in
the primitive
marked top-cell
readouts,
\begin{equation}\label{eq:dyadic-f134-trace-restriction-three}
\boxed{
\begin{aligned}
\operatorname{Tr}_{F,\vartheta}
\circ\operatorname{res}_{F/\mathbf Q_2}
&=
3\operatorname{Tr}_{\Gamma,\vartheta},\\
(\operatorname{Tr}_{F,\bf1}\bmod2)
\circ\operatorname{res}_{F/\mathbf Q_2}
&=
3(\operatorname{Tr}_{\Gamma,\bf1}\bmod2)
=
\operatorname{Tr}_{\Gamma,\bf1}\bmod2.
\end{aligned}}
\end{equation}
In particular the
top cyclotomic
restriction
is an \emph{integral
isomorphism}
because $3$ is
a unit at $2$.
After the usual
local fundamental
class normalization,
the same equation
is the classical
Tate invariant
formula
\begin{equation}\label{eq:dyadic-f134-local-tate-trace-normalization}
\operatorname{inv}_{F}
(\operatorname{res}_{F/\mathbf Q_2}x)
=
3\operatorname{inv}_{\mathbf Q_2}(x),
\qquad
\operatorname{inv}_{\mathbf Q_2}
(\operatorname{cor}_{F/\mathbf Q_2}y)
=
\operatorname{inv}_{F}(y).
\end{equation}
The marked top-cell
readouts are normalized
by the chosen
relative relation
faces; identifying
their generators
with a given
arithmetically
normalized
fundamental class
is still a choice
of a $2$-adic
unit sign/scale.

For precisely these
two rank-one
coefficient modules,
the \emph{formal}
full marked
two-face complex
has the same
cohomology ledger
as the intrinsic
local Galois
cochain complex:
\begin{equation}\label{eq:dyadic-f134-global-rankone-ledgers}
\begin{array}{c|ccc}
&H^0&H^1&H^2\\ \hline
\mathcal J_\Gamma(\mathbf Z_2)
&\mathbf Z_2&\mathbf Z_2^2&\mathbf Z/2\\
\mathcal J_\Gamma(\mathbf Z_2(1))
&0&\mathbf Z_2^2\oplus\mathbf Z/2&
\mathbf Z_2 .
\end{array}
\end{equation}
The statement is
\emph{coefficient-specific}:
it does not assert
that the naked
two-relator marked
complex is an
exact projective
resolution for
all continuous
$\Gamma$-modules.
\end{theorem}

\begin{proof}
For a unit-valued
cyclotomic character,
every continuous
crossed derivation
to $\mathbf Z_2(1)$
annihilates the
pro-odd part of
any procyclic
source element.
More generally,
this holds for
any rank-one
$\mathbf Z_2$
module whose
action has
pro-$2$ image:
on the pro-odd
part the action
is trivial,
and a continuous
homomorphism
from a pro-odd
group to the
additive pro-$2$
group is zero.
Consequently
\begin{equation}\label{eq:dyadic-f134-omega2-crossed-derivation}
\boxed{\qquad
D(g^{\omega_2})=D(g)
\qquad}
\end{equation}
when evaluating
the completed
marked Fox
derivatives in
either of these
two coefficient
modules.
This \emph{does
not} identify
the actual
group words
$g^{\omega_2}$
and $g$;
it is only
a statement
about the
specialized
crossed derivation.

For trivial
coefficients,
evaluating the
literal tame
and wild
Roe--Turturean
relations yields
the rows
$(0,-1,0,0)$
and $(0,3,2,0)$.
For the
cyclotomic
character, use
\[
\vartheta(\sigma)
=\vartheta(\tau)=1,\quad
\vartheta(x_0)=-1,\quad
\vartheta(x_1)=-1/3
\]
and the same
specialized
$\omega_2$
derivation lemma.
Substitution in
the auxiliary
wild source words
gives the rows
$(0,-1,0,0)$
and $(0,-3,0,0)$.
The respective
principal degree-zero
coboundaries
are
$d^0=(\vartheta(g)-1)_g$,
proving
\eqref{eq:dyadic-f134-full-marked-two-twist-rows}.

For trivial
coefficients
the tame row
contains a unit,
and the
remaining wild
row has one
invariant factor
$2$.  Its
cokernel is
$\mathbf Z/2$,
detected by
$w+3t$ mod two.
For twisted
coefficients
the wild row
is exactly
three times
the tame row,
so the cokernel
is freely
generated by
$w-3t$.
The integral
degree-zero
map for the
twist is twice
a primitive
vector; the
usual elementary
divisor argument
therefore gives
the ledgers
\eqref{eq:dyadic-f134-global-rankone-ledgers}.
They agree with
local class
field theory,
integral Kummer
theory, and Tate
duality for
$\mathbf Q_2$.

The three-sheet
covering has
one copy of
the genuine
marked wild
relation face
above each
coset representative
$t_i=\tau^i$.
Both coefficient
actions send
$\tau$ to one.
Therefore a
top cochain
whose tame
coordinate
has been killed
by the unit
tame relation
pivot and whose
wild coordinate
is $c$ pulls
back to equal
values $c$
on all three
lifted wild
faces.
The three
tame-face
cancellations
and the two
primitive wild
Tietze
cancellations
preserve their
total face
coordinate,
by the relation
product
$R_0R_1R_2$
in fixed133.
The resulting
top trace is
$3c$, proving
\eqref{eq:dyadic-f134-trace-restriction-three}
on the formal
face cohomology
quotients.

For the actual
Galois groups,
the canonical
cohomological
restriction and
corestriction
satisfy
$\operatorname{cor}
\operatorname{res}=3$
because $[F:\mathbf Q_2]=3$.
The local
fundamental
class comparison
is normalized
by the standard
Tate invariants,
which obey
\eqref{eq:dyadic-f134-local-tate-trace-normalization}
\cite{NSW}.
The marked
cover comparison
of fixed133
is an actual
Fox resolution
of $H(2)$,
so its
top-face
readout is
primitive and
nonzero.
The explicit
restriction of
the global
wild face
therefore has
nonzero primitive
image and is
an isomorphism
after inverting
the odd factor
$3$.
This reconciles
the marked
face normalization
with the
intrinsic
local Galois
invariant.
No claim of
a universally
exact full-$\Gamma$
group-algebra
two-relator
resolution has
been made.
\end{proof}

\begin{corollary}[Fox--Tate finite certificate without inverse Nielsen words]
\label{cor:dyadic-f134-primitive-trace-certificate}
For the selected
index-three
$S_3$ Sylow
preimage
$H=G_{\mathbf Q_2(\sqrt[3]2)}$,
the following
finite marked
source data
determine the
two integral
top cohomology
readouts,
the ordinary
and cyclotomic
Bockstein
placement, and
the index-three
restriction
factor:
\begin{equation}\label{eq:dyadic-f134-final-trace-certificate}
\boxed{
\begin{gathered}
(d_{\bf1}^0,d_{\bf1}^1)
=(0,J_{\bf1}),\\
(d_{\vartheta}^0,d_{\vartheta}^1)
=\bigl(
(0,-2,-2,-2,-4/3,-4/3,-4/3)^t,
J_{\vartheta}\bigr),\\
\operatorname{Tr}_F=(1,1,1),
\qquad
\operatorname{Tr}_F\operatorname{res}
=3\operatorname{Tr}_{\mathbf Q_2}.
\end{gathered}}
\end{equation}
All displayed
matrices are
integral over
$\mathbf Z_2$;
the only
denominators are
odd units.
The certificate
recovers
$H^2(H,\mathbf Z_2)
=\mathbf F_2$
and
$H^2(H,\mathbf Z_2(1))
=\mathbf Z_2$
\emph{directly}
from the three
Schreier relation
faces, without
computing the
full nonlinear
Nielsen inverse
$\alpha^{-1}$.
It complements,
rather than
replaces, the
group-level
Schreier
syzygy closure
of fixed133.
\end{corollary}

\begin{proof}
Theorem~\ref{thm:dyadic-f134-exact-two-fox-matrices}
gives the two
finite integral
Fox rows and their
top-coboundary
images.
Theorem~\ref{thm:dyadic-f134-top-trace-bockstein-shift}
gives the primitive
relation-face
trace and both
Bockstein placements.
Theorem~\ref{thm:dyadic-f134-top-transfer-three}
supplies the
index-three
compatibility.
All source data
are finite
integral matrices
and marked
coset words;
the inverse
Nielsen relation
is not required
for these
particular
cohomological
readouts.
\end{proof}

\begin{warning}[Scope of the relative Fox--Tate trace]
\label{warn:dyadic-f134-rankone-trace-scope}
The global
two-relation
marked complex
is used
\emph{only}
for the two
rank-one
coefficient
systems
$\mathbf Z_2$
and
$\mathbf Z_2(1)$.
On the
index-three
pro-$2$
Schreier cover,
the fixed133
seven-generator
complex is a
genuine
group-cohomological
Fox resolution;
its integral
top trace is
therefore exact.
For a general
higher-rank
Galois family
one cannot
identify the
complete
full-group
marked face
cokernel with
$H^2$ without
checking the
appropriate
projective
syzygy
conditions.
Nor does the
sum-of-faces
operation
make a
non-multiplicative
corestriction
into a
strict
cup-algebra
homomorphism.
The present
certificate
closes the
two specified
rank-one
Fox--Tate
readouts
and their
odd-index
restriction
normalization,
not the
global
presentation-independent
Fox gauge.
\end{warning}
